\documentclass[11pt,a4paper]{article}

\usepackage[margin=0.8in]{geometry}
\usepackage[T1]{fontenc}
\usepackage[utf8]{inputenc}
\usepackage{lmodern}
\usepackage{microtype}
\usepackage[svgnames]{xcolor}
\usepackage[hidelinks]{hyperref}
\usepackage{xurl}
\usepackage{enumitem}
\usepackage{titlesec}
\usepackage{tabularx}
\usepackage{array}
\usepackage[
    backend=biber,
    authordate,
    sorting=ydnt,
    maxbibnames=99,
    dashed=false,
    giveninits=true,
    uniquename=init,
    date=year,
    doi=false,
    isbn=false,
    url=false,
    eprint=false
]{biblatex-chicago}

\addbibresource{publications.bib}

\definecolor{accent}{HTML}{1F4E79}
\definecolor{textgray}{HTML}{4F5B66}
\definecolor{softline}{HTML}{D9E2EC}

\pagestyle{empty}
\setlength{\parindent}{0pt}
\setlength{\parskip}{0pt}
\setlength{\emergencystretch}{2em}
\setlist[itemize]{leftmargin=1.3em, topsep=3pt, itemsep=2pt, parsep=0pt, partopsep=0pt}
\urlstyle{same}

\titleformat{\section}[block]{\Large\bfseries\color{accent}}{}{0em}{\strut}[\vspace{0.15em}\color{softline}\titlerule]
\titlespacing*{\section}{0pt}{1.0em}{0.8em}

\newcolumntype{Y}{>{\raggedright\arraybackslash}X}

\DeclareFieldFormat[article,inproceedings]{title}{\mkbibemph{#1}}
\renewcommand*{\bibfont}{\small}
\setlength{\bibitemsep}{0.45\baselineskip}
\AtEveryBibitem{%
    \clearfield{month}%
    \clearfield{number}%
    \clearlist{language}%
}

\newcommand{\cvsubsection}[1]{%
    \vspace{0.15em}
    {\normalsize\bfseries\color{accent!90}#1}\par
    \vspace{0.2em}
}

\newcommand{\cventry}[4]{%
    \begin{tabularx}{\textwidth}{@{}Y>{\raggedleft\arraybackslash}p{0.28\textwidth}@{}}
        \textbf{#1} & \textcolor{textgray}{#2} \\
        \textit{#3} &
    \end{tabularx}
    #4\par\vspace{0.45em}
}

\newcommand{\projectentry}[3]{%
    \begin{tabular*}{\textwidth}{@{\extracolsep{\fill}}lr@{}}
        \textbf{#1} & {\small\href{#2}{\nolinkurl{#3}}} \\
    \end{tabular*}
}

\begin{document}

\begin{center}
    {\fontsize{24}{28}\selectfont\bfseries Kaihuang Chen}\par
    \vspace{4pt}
    {\color{textgray}\normalsize Ph.D. in Applied Mathematics}\par
    {\color{textgray}\normalsize Large-Scale Optimization \textbar{} GPU-Accelerated Scientific Computing}\par
    \vspace{6pt}
    {\small Hong Kong SAR, China }
\end{center}

\vspace{0.35em}
{\color{softline}\hrule height 0.9pt}

\section*{Profile}
Completed doctoral research in large-scale optimization and high-performance scientific computing, with emphasis on algorithm design and GPU implementation for linear and quadratic programming. Research focuses on operator splitting methods, scalable sparse computation, and reproducible numerical experiments for large optimization models.

\section*{Research Interests}
Numerical optimization; large-scale convex optimization; high-performance computing; GPU-accelerated optimization algorithms; linear programming; quadratic programming; operator splitting methods; and sparse scientific computing.

\section*{Education}
\cventry{The Hong Kong Polytechnic University (PolyU)}{Sep 2020 -- Feb 2026}{Department of Applied Mathematics}{Ph.D. in Applied Mathematics}
\cventry{Hunan University}{Sep 2016 -- Jun 2020}{School of Mathematics}{B.Sc. in Information and Computing Science}

\section*{Publications}
\nocite{*}
\cvsubsection{Journal Article}
\printbibliography[heading=none,type=article,notkeyword=preprint]

\cvsubsection{Conference Publication}
\printbibliography[heading=none,type=inproceedings]

\cvsubsection{Preprints}
\printbibliography[heading=none,keyword=preprint]

\section*{Software Projects}
\projectentry{HPR-LP}{https://github.com/PolyU-IOR/HPR-LP}{github.com/PolyU-IOR/HPR-LP}
\begin{itemize}
    \item Julia/GPU solver for large-scale linear programming.
    \item Developed for reproducible large-scale benchmarking and high-performance optimization experiments.
\end{itemize}

\projectentry{HPR-QP}{https://github.com/PolyU-IOR/HPR-QP}{github.com/PolyU-IOR/HPR-QP}
\begin{itemize}
    \item Julia/CUDA solver for large-scale convex composite quadratic programming.
    \item Implements a dual Halpern--Peaceman--Rachford framework with practical enhancements for large sparse instances.
    \item Designed for GPU-accelerated research experiments and scalable numerical testing.
\end{itemize}

\section*{Technical Skills}
\begin{tabularx}{\textwidth}{@{}>{\bfseries}p{0.18\textwidth}Y@{}}
Programming & Julia, MATLAB, CUDA, C/C++ \\
Research Areas & Numerical optimization, convex optimization, LP, QP, operator splitting, scientific computing, sparse linear algebra, GPU computing \\
\end{tabularx}

\end{document}